\section{LMI conditions and numerical settings}\label{app:lmi}

\paragraph{Collocated model}
Writing Eq.~\eqref{M-eq:dev} in the state $\xi$ of Section~\ref{M-sec:cert-model} gives Eq.~\eqref{M-eq:coll} with
\begin{equation}\label{eq:app-A0}
A_0=\begin{bmatrix}
0&1&0&0&\bm 0\\
-k/m&-c/m&-K_a\omega_a^2/m&-2\zeta_a\omega_aK_a/m&\bm 0\\
0&0&0&1&\bm 0\\
0&0&-\omega_a^2&-2\zeta_a\omega_a&\bm 0\\
\hat{\bm d}_0&\bm 0&\bm 0&\bm 0&\hat D
\end{bmatrix},
\end{equation}
$B=[0,\ K_a/m,\ 0,\ 1,\ \bm 0]^\top$, where $\hat{\bm d}_0\in\mathbb R^N$ and $\hat D\in\mathbb R^{N\times N}$ are column 0 and columns $1,\dots,N$ of rows $1,\dots,N$ of $(2/T)D_N$, $D_N$ being the Chebyshev differentiation matrix on the points $\cos(\pi j/N)$, $j=0,\dots,N$ \citep{trefethen2000spectral}: the last $N$ rows state $\tfrac{\mathrm d}{\mathrm dt}\tilde x(t+\vartheta_j)=\tfrac2T\sum_{l=0}^{N}(D_N)_{jl}\,\tilde x(t+\vartheta_l)$. The last of these rows, $d_N^\top=e_{4+N}^\top A_0$, is the derivative of the interpolant at $\vartheta=-T$ used in Section~\ref{supp:delayed-model}. Table~\ref{tab:nconv} gives the certified kick for $N$ and $2N$ nodes, and for $1.5N$ at some points, $N$ being the default of Section~\ref{M-sec:cert-model}, which reproduces the stored C2 certificates.
% src: src/c2_certificate.py (cheb, model, history_nodes); src/c3_periodic.py (model); src/c4_comb.py (model: x'(t - T) from row N of (2/T) D_N); verification/paper/certs_nconv_summary.py/certs_nconv.json (merged from certs_nconv_a/b/c.json, certs_nconv_extra.json and revision_runs.json beta_vs_N); values at the default N equal results/c2_validate.npz cert at all 11 C2 points

\begin{table}[!tbp]
\centering
\caption{Certified kick $\beta$ (\si{\milli\metre\per\second}) of C2 against the number $N$ of collocation nodes, with the command bound of the original C2 solves (before the rescaling to $\bar b$); --: not computed; ${}^{\dagger}$no candidate passed the independent check (with 24 and 36 nodes, $\beta=1.133$). DVF\,+\,D: $g_1=500$, $g_2=\SI{4000}{\volt\second\per\metre}$ (Section~\ref{M-sec:delayed}).}
\label{tab:nconv}
\small
\begin{tabular}{@{}llccccr@{}}
\toprule
Model & Point & $N$ & $\beta$ at $N$ & at $1.5N$ & at $2N$ & Change at $2N$\\
\midrule
C2 & \SI{1090}{\rpm}, $0.5\alim$ & 24 & 8.93 & -- & 8.49 & $-5.0$\,\%\\
C2 & \SI{1090}{\rpm}, $0.9\alim$ & 24 & 6.20 & 6.14 & 6.10 & $-1.6$\,\%\\
C2 & \SI{1090}{\rpm}, $0.98\alim$ & 24 & 4.85 & -- & 4.76 & $-1.7$\,\%\\
C2 & \SI{1420}{\rpm}, $0.5\alim$ & 24 & 8.31 & -- & 8.24 & $-0.8$\,\%\\
C2 & \SI{1420}{\rpm}, $0.98\alim$ & 24 & 2.08 & -- & 2.08 & $-0.05$\,\%\\
C2 & \SI{2500}{\rpm}, $0.5\alim$ & 20 & 2.59 & -- & 2.59 & $-0.06$\,\%\\
C2 & \SI{2500}{\rpm}, $0.9\alim$ & 20 & 3.06 & 3.06 & 3.06 & $0.0$\,\%\\
C2 & \SI{2500}{\rpm}, $0.98\alim$ & 20 & 1.14 & 1.13 & ${}^{\dagger}$ & --\\
C2 & \SI{2870}{\rpm}, $0.3\alim$ & 20 & 8.74 & 8.61 & 8.61 & $-1.5$\,\%\\
C2 & \SI{2870}{\rpm}, $0.5\alim$ & 20 & 5.51 & -- & 5.55 & $+0.6$\,\%\\
C2 & \SI{2870}{\rpm}, $0.98\alim$ & 20 & 4.83 & -- & 4.82 & $-0.2$\,\%\\
DVF\,+\,D & \SI{2500}{\rpm}, \SI{22}{\milli\metre} & 20 & 3.56 & 3.54 & 3.53 & $-0.9$\,\%\\
\bottomrule
\end{tabular}
\end{table}
% src: verification/paper/certs_nconv.json (points: 1090_0.5 8.9287/8.4856; 1090_0.9 6.1961/6.1447/6.0975; 1090_0.98 4.8459/4.7639; 1420_0.5 8.3122/8.2416; 1420_0.98 2.0828/2.0818; 2500_0.5 2.5869/2.5853; 2500_0.9 3.0579/3.0568/3.0578; 2500_0.98 1.1397 (N 20), 1.1331 (24), 1.1329 (30), 1.1327 (36), check failed at 40; 2870_0.3 8.7429/8.6089/8.6076; 2870_0.5 5.5129/5.5478; 2870_0.98 4.8266/4.8169); verification/paper/revision_runs.json (beta_vs_N nconv_comb: 3.5616/3.5413/3.5296 mm/s at N = 20/30/40)

\paragraph{Conditions of C2}
The program maximises $\beta$ subject to Eqs.~\eqref{M-eq:lmi-dec}--\eqref{M-eq:lmi-size}, with the left side of Eq.~\eqref{M-eq:lmi-dec} reduced by $\varepsilon I$, $S\ge\varepsilon$, $\varepsilon=10^{-9}$, and $\beta\ge0$. It has $n_\xi(n_\xi+1)/2+n_\xi+2$ unknowns, 326--704 for $n_\xi=4+N=24$--36, in four LMI blocks of order at most $n_\xi+1$. The decay margin is $\lambda=0.05\lambda_0$, $\lambda_0=-\max\operatorname{Re}\operatorname{eig}\mathcal A$ with $aH=\overline{aH}$, and $0.2\lambda_0$ when no candidate passed the check. The first solve uses coordinates scaled by the peak response of each state to a \SI{1}{\milli\metre\per\second} kick, the next two coordinates whitened by the previous $W$ (re-centring); at all 48 points of Section~\ref{M-sec:cert-c2} the kept candidate came from a re-centred solve, five of them with the solver status `almost solved', which the check does not rely on.
% src: src/c2_certificate.py (lmi: eps 1e-9, blocks M1-M4; LAM_FRAC (0.05, 0.2); scaling; certify: N_PASS 3); results/c2_validate.npz (msg: pass 1 or 2 at all 48 points; AlmostSolved at 5); unknowns n_xi(n_xi+1)/2+n_xi+2 evaluated for n_xi = 24 and 36

\paragraph{Periodic certificate (C3)}
The tooth period is divided into $M$ segments of length $\tau_s=T/M$, with nodes $t_k=k\tau_s$. On segment $k$, with $s=(t-t_k)/\tau_s\in[0,1]$,
\begin{equation}\label{eq:Wt}
W(t)=(1-s)W_k+sW_{k+1},\qquad Y(t)=(1-s)Y_k+sY_{k+1},\qquad S(t)=S_k,
\end{equation}
with $W_M=W_0$ and $Y_M=Y_0$, so that $\mathcal V(t,\xi)=\xi^\top W(t)^{-1}\xi$ is continuous and $T$-periodic. As $\mathrm d(W^{-1})/\mathrm dt=-W^{-1}\dot WW^{-1}$, the decrease condition becomes
\begin{equation}\label{eq:lmi-per}
\begin{bmatrix}\Delta_kW-\mathcal A(t)W(t)-W(t)\mathcal A(t)^\top-2\lambda W(t) & BS_k-Y(t)^\top\\ S_kB^\top-Y(t) & 2S_k\end{bmatrix}\succ0,
\end{equation}
with $\Delta_kW=(W_{k+1}-W_k)/\tau_s$, the derivative of $W(t)$ on segment $k$. On each segment, $aH(t)$ is enclosed by a line and a residual bound, and the sector condition uses the segment level $u_{0,k}=\Vmax-\bar b_k$, with $\bar b_k\ge\lvert b(t)\rvert$ there (see \emph{Enclosure and command bound} below).
% src: README §14.1, §14.2; src/c3_periodic.py (aH, orbit, enclosure); |b| range re-evaluated in closed form with c3_periodic.orbit (0.031-7.870 V at 2500 rpm, 0.98 a_lim)

\begin{proposition}[Periodic model, C3]\label{prop:c3}
Let $\mathcal A(t)$ be $T$-periodic and, on segment $k$, $aH(t)=c_{0,k}+c_{1,k}s+r(t)$ with $\lvert r(t)\rvert\le\varepsilon_k$, and $u_{0,k}>0$. Let $W$, $Y$ and $S$ follow Eq.~\eqref{eq:Wt} with $W_k\succ0$. If, for every $k$, the matrix of Eq.~\eqref{eq:lmi-per} with $\lambda=0$, the enclosure in place of $aH(t)$ and $\sigma$ in place of $s^2$, is positive definite at the six vertices $(s,\sigma)\in\{(0,0),(\tfrac12,0),(1,1)\}$, $r=\pm\varepsilon_k$, and Eqs.~\eqref{M-eq:lmi-incl} and~\eqref{M-eq:lmi-valid} hold at every node with $W_k$, $Y_k$ and the level $\min(u_{0,k-1},u_{0,k})$ (indices modulo $M$), then the conclusions of Proposition~\ref{M-prop:c2} hold for the tube $\{(t,\xi):\xi^\top W(t)^{-1}\xi\le1\}$. If Eq.~\eqref{M-eq:lmi-size} holds at both nodes of a segment, $\beta e\in\mathcal E\big(W(t)^{-1}\big)$ on that segment.
\end{proposition}
The proof, which reduces the periodic decrease condition on each segment to its six vertices by convexity, is in Section~\ref{supp:proof-c3}.

\paragraph{Conditions of C3}
On each segment the decrease condition~\eqref{eq:lmi-per} is imposed at the six vertices of Proposition~\ref{prop:c3}, the two vertices $r=\pm\varepsilon_k$ of each $(s,\sigma)$ being merged into one LMI by Petersen's lemma, exact for a scalar uncertainty \citep{petersen1987stabilization} (Section~\ref{supp:cert-vertex}). With $M=16$ there are $M[n_\xi(n_\xi+1)/2+n_\xi+1]+1$ unknowns and $3M$ multipliers, 7009 scalars at $N=24$. The decay margin is $0.05$, 0.02 or 0.01 times the Floquet decay rate, the next being tried when no candidate passes. The single solve starts in the coordinates of the C2 certificate; a second, re-centred pass added only 0.4\,\% at \SI{2500}{\rpm} and $0.9\alim$. Merging made the solves 3--4 times faster; the merged problem needed balanced multiplier blocks to start, and with the fallbacks 47 points were certified, the one at \SI{2870}{\rpm} and $0.98\alim$ only with separate vertices.
% src: src/c3_periodic.py (corners, lmi: merged block, P_SCALE 10; LAM_FRAC (0.05, 0.02, 0.01)); README §14.2, §14.9; unknowns 16 x 435 + 1 + 48 = 7009 for n = 28

\paragraph{Enclosure and command bound}
The line $c_{0,k}+c_{1,k}s$ is a least-squares fit to $aH$ at 2048 points of segment $k$, and $\varepsilon_k$ bounds the residual, including its variation between the points (Section~\ref{supp:cert-vertex}). The level $u_{0,k}$ uses the largest of 257 samples of $\lvert b(t)\rvert$ on the segment, plus $\max\lvert\ddot b\rvert\,\Delta t^2/8$ ($\Delta t$ the sample spacing), plus $g$ times a bound on the forced velocity carried by the harmonics above the 40th (summed to the 20\,000th and bounded analytically beyond), with the slice-force harmonics in closed form; $\bar b$ is the largest of these segment bounds. C2 was first solved with the largest of 4096 samples plus the same inter-sample bound, for 40 harmonics whose slice-force coefficients came from a 4096-point discrete Fourier transform with the entry jump on a sample. $\bar b$ exceeds that bound by 0.4--\SI{2.1}{\milli\volt} at the 48 points and by 0.1--\SI{3.3}{\milli\volt} on the surface grid. The forced command of the simulator, computed with 30 axial layers, exceeds the closed form by up to \SI{8.7}{\percent} (at most \SI{0.33}{\volt}) at the operating points of Section~\ref{M-sec:feasibility}, so the simulator's margin is slightly narrower than the certificates assume; every certificate nevertheless lies below the simulated basin. As Eq.~\eqref{M-eq:lmi-dec} is homogeneous in $(W,Y,S)$ and the quadratic forms of Eqs.~\eqref{M-eq:lmi-incl}--\eqref{M-eq:lmi-size} scale with them, a certificate at the level $u_0$ remains valid at $u_0-d$ after scaling $(W,Y,S)$ by $\big((u_0-d)/u_0\big)^2$, which scales $\beta$ by $(u_0-d)/u_0$; all C2 results are rescaled in this way to $\bar b$.
% src: src/c3_periodic.py (enclosure: N_ENC 2048, aH_rate_bound, orbit, velocity_tail, lmi); src/c2_certificate.py (orbit); src/c3_summary.py (margin_errata); results/c3_summary.json (errata: extra_mV 0.414-2.079 and 0.131-3.275); README §14.2, §14.9; README §3 (cf_vs_td_layers: up to +8.7 %); results/feas_g1500.npz (cert_demand_td - cert_demand_cf at most 0.98 N = 0.33 V at the 129 settled points with demand >= 1 N)

\paragraph{Independent check}
A candidate passes if $W\succ0$ with a condition number below $10^{10}$ and if the smallest computed eigenvalue of every decrease matrix at $\lambda=0$ exceeds $10^3(n_\xi+1)\epsilon_{\mathrm{mach}}$ times its largest absolute eigenvalue. This margin lies well above the backward error of the symmetric eigensolver, but the test is made in floating point. It is not an interval-arithmetic or verified-Cholesky proof \citep{rump2006verification}, and for C2, whose matrices are formed in the re-centred coordinates, the tolerance does not cover the rounding made in forming them. For C3 the test is applied at all $6M$ vertices separately, also when they were merged in the solve, with a scale that also covers the rounding in $\Delta_kW$ and in $\mathcal AW$ (Section~\ref{supp:cert-vertex}). Inclusion, validity and the size $\beta=(e^\top W^{-1}e)^{-1/2}$ are evaluated at every node through Cholesky factors, and an exceeded bound is repaired by one common scaling of all $(W_k,Y_k,S_k)$, which keeps every decrease LMI by homogeneity. The largest candidate that passed is kept, and the smaller of its checked size and the solver's $\beta$ is reported.
% src: src/c2_certificate.py (check: condition number of W below 1e10, tolerance 1e3 (n+1) eps max|lambda|; certify); src/c3_periodic.py (check)

\begin{table}[!tbp]
\centering
\caption{Numerical settings and cost of the certificates; $\lambda_0$ is the linear decay rate (C2) or the Floquet decay rate (C3).}
\label{tab:lmi-settings}
\small
\begin{tabular}{@{}>{\raggedright\arraybackslash}p{0.17\textwidth}>{\raggedright\arraybackslash}p{0.38\textwidth}>{\raggedright\arraybackslash}p{0.38\textwidth}@{}}
\toprule
 & C2 (averaged model) & C3 (periodic model)\\
\midrule
Solver & Clarabel; gap (absolute, relative) and feasibility tolerances $10^{-10}$, at most 400 iterations (48 points, step entry); defaults (surface) & Clarabel defaults, at most 120 iterations, direct solver faer, one thread\\
Fallbacks & $\lambda=0.2\lambda_0$ if no candidate passes at $0.05\lambda_0$ (used for 4 of the 287 surface certificates) & $\lambda=0.05$, 0.02, $0.01\lambda_0$; static regularisation $10^{-7}$ instead of $10^{-8}$; separate vertices\\
Strictness & $\varepsilon=10^{-9}$ & $\varepsilon=10^{-9}$\\
Discretisation & $N=16+4\lceil T\omega_n/2\pi\rceil=20$--32 & as C2; $M=16$ (8 and 32 in Section~\ref{supp:cert-decomp})\\
Solves & three: scaled, then two re-centred & one, in the coordinates of the C2 certificate\\
Size & 326--704 unknowns & 7009 scalar unknowns at $N=24$, multipliers included\\
Time, memory & about \SI{2}{\hour} for the campaign, including its simulations, on four cores & 4--\SI{10}{\minute} per certificate, 2--5\,GB (4.7\,GB with separate vertices at $N=24$; about twice with $M=32$, 5.2\,GB at $N=20$); about \SI{8}{\hour} for the campaign on four cores\\
\bottomrule
\end{tabular}
\end{table}
% src: src/c2_certificate.py (TIGHT, FAST, LAM_FRAC, N_PASS, nodes_for); src/c3_periodic.py (SOLVER, REG_RETRY, ATTEMPTS, LAM_FRAC, N_PASS, M_C); results/c2_surface.npz (msg: lam 0.2 sigma at 4 points); README §10, §14.6, §14.9

\paragraph{Settings and model runs}
Table~\ref{tab:lmi-settings} lists the settings; sizes reproduce to about 0.1\,\%. Critical kicks of the collocated models used the true dead zone, held over each of 256 steps per tooth period, with the linear part integrated exactly (C2) or by fourth-order Magnus propagators (C3) \citep{blanes2009magnus}; a run diverged if its state norm was not finite or exceeded $10\Delta v$ after 1200 tooth periods. Logarithmic bisection on $[10^{-4},3]$\,m/s used nine steps (about 2\,\%) and the minimum over four phases in C2, and 14 steps (0.06\,\%) at each of the 16 phases in C3; the time discretisation moves the result by about 0.1--0.2\,\%. The off-axis runs applied boundary points of the certified ellipsoid at four phases, at revolution 100 of a face entry with the \SI{9}{\volt} clip, and classified the following 300 revolutions as in Section~\ref{M-sec:bistability}. The C2 runs used the full surface memory, the C3 runs the last four passes, as in the original campaign, their starts lying inside the validity bound; four passes overstate the chip once the tool has left the cut for longer, which tends to favour chatter, and none of these runs chattered.
% src: src/c2_certificate.py (run_model, model_critical_kick, td_offaxis); src/c3_periodic.py (propagators, step_maps, model_kick_by_phase, NBIS, L_SIM, td_offaxis); README §14.9 (reproducibility 0.1 %, time step 0.1-0.2 %)
